\chapter{2024年北京卷}

\section{单选题}

\begin{problem}
已知集合 \(M=\{x\mid -3<x<1\}\)，\(N=\{x\mid -1\le x<4\}\)，则 \(M\cup N=\)（\quad）

\choices
    {\(\{x\mid -1\le x<1\}\)}
    {\(\{x\mid x>-3\}\)}
    {\(\{x\mid -3<x<4\}\)}
    {\(\{x\mid x<4\}\)}
\end{problem}

\begin{problem}
已知 \(\frac{z}{\i}=-1-\i\)，则 \(z=\)（\quad）

\choices
    {\(-1-\i\)}
    {\(-1+\i\)}
    {\(1-\i\)}
    {\(1+\i\)}
\end{problem}

\begin{problem}
圆 \(x^2+y^2-2x+6y=0\) 的圆心到直线 \(x-y+2=0\) 的距离为（\quad）

\choices
    {\(\sqrt{2}\)}
    {2}
    {3}
    {\(3\sqrt{2}\)}
\end{problem}

\begin{problem}
在 \(\left( x-\sqrt{x} \right)^4\) 的展开式中，\(x^3\) 的系数为（\quad）

\choices
    {6}
    {\(-6\)}
    {12}
    {\(-12\)}
\end{problem}

\begin{problem}
设 \(\bs{a}\)，\(\bs{b}\) 是向量，则“\(\left( \bs{a}+\bs{b} \right)\cdot\left( \bs{a}-\bs{b} \right)=0\)”是“\(\bs{a}=-\bs{b}\) 或 \(\bs{a}=\bs{b}\)”的（\quad）

\choices
    {充分不必要条件}
    {必要不充分条件}
    {充要条件}
    {既不充分也不必要条件}
\end{problem}

\begin{problem}
设函数 \(f(x)=\sin \omega x\)（\(\omega>0\)）. 已知 \(f(x_1)=-1\)，\(f(x_2)=1\)，且 \(\lvert x_1-x_2\rvert\) 的最小值为 \(\frac{\pi}{2}\)，则 \(\omega=\)（\quad）

\choices
    {1}
    {2}
    {3}
    {4}
\end{problem}

\begin{problem}
生物丰富度指数 \(d=\frac{S-1}{\ln N}\) 是河流水质的一个评价指标，其中 \(S\)，\(N\) 分别表示河流中的生物种类数与生物个体总数. 生物丰富度指数 \(d\) 越大，水质越好. 如果某河流治理前后的生物种类数 \(S\) 没有变化，生物个体总数由 \(N_1\) 变为 \(N_2\)，生物丰富度指数由 \(2.1\) 提高到 \(3.15\)，则（\quad）

\choices
    {\(3N_2=2N_1\)}
    {\(2N_2=3N_1\)}
    {\(N_2^2=N_1^3\)}
    {\(N_2^3=N_1^2\)}
\end{problem}

\begin{problem}
如图，在四棱锥 \(P-ABCD\) 中，底面 \(ABCD\) 是边长为 4 的正方形，\(PA=PB=4\)，\(PC=PD=2\sqrt{2}\)，该棱锥的高为（\quad）

\bitmapfigure[max width=5.2cm]{img/2024/beijing/q8_fig1.png}

\choices
    {1}
    {2}
    {\(\sqrt{2}\)}
    {\(\sqrt{3}\)}
\end{problem}

\begin{problem}
已知 \((x_1,y_1)\)，\((x_2,y_2)\) 是函数 \(y=2^x\) 的图象上两个不同的点，则（\quad）

\choices
    {\(\log_2 \frac{y_1+y_2}{2}<\frac{x_1+x_2}{2}\)}
    {\(\log_2 \frac{y_1+y_2}{2}>\frac{x_1+x_2}{2}\)}
    {\(\log_2 \frac{y_1+y_2}{2}<x_1+x_2\)}
    {\(\log_2 \frac{y_1+y_2}{2}>x_1+x_2\)}
\end{problem}

\begin{problem}
已知 \(M=\left\{ (x,y)\mid y=x+t\left( x^2-x \right),\,1\le x\le 2,\,0\le t\le 1 \right\}\) 是平面直角坐标系中的点集. 设 \(d\) 是 \(M\) 中两点间距离的最大值，\(S\) 是 \(M\) 表示的图形的面积，则（\quad）

\choices
    {\(d=3\)，\(S<1\)}
    {\(d=3\)，\(S>1\)}
    {\(d=\sqrt{10}\)，\(S<1\)}
    {\(d=\sqrt{10}\)，\(S>1\)}
\end{problem}

\section{填空题}

\begin{problem}
抛物线 \(y^2=16x\) 的焦点坐标为\fillinblank{}.
\end{problem}

\begin{problem}
在平面直角坐标系 \(xOy\) 中，角 \(\alpha\) 与角 \(\beta\) 均以 \(Ox\) 为始边，它们的终边关于原点对称. 若 \(\alpha\in \left[ \frac{\pi}{6},\frac{\pi}{3} \right]\)，则 \(\cos \beta\) 的最大值为\fillinblank{}.
\end{problem}

\begin{problem}
若直线 \(y=k(x-3)\) 与双曲线 \(\frac{x^2}{4}-y^2=1\) 只有一个公共点，则 \(k\) 的一个取值为\fillinblank{}.
\end{problem}

\begin{problem}
汉代刘歆设计的“铜嘉量”是龠，合，升，斗，斛五量合一的标准量器，其中升量器，斗量器，斛量器的形状均可视为圆柱. 若升，斗，斛量器的容积成公比为 10 的等比数列，底面直径依次为 \(65\mathrm{~mm}\)，\(325\mathrm{~mm}\)，\(325\mathrm{~mm}\)，且斛量器的高为 \(230\mathrm{~mm}\)，则斗量器的高为\fillinblank{}\(\mathrm{~mm}\)，升量器的高为\fillinblank{}\(\mathrm{~mm}\).（不计量器的厚度）
\end{problem}

\begin{problem}
设 \(\{a_n\}\) 与 \(\{b_n\}\) 是两个不同的无穷数列，且都不是常数列. 记集合 \(M=\{k\mid a_k=b_k,\ k\in \N^*\}\)，给出下列 4 个结论：

\begin{circlelist}
\item 若 \(\{a_n\}\) 与 \(\{b_n\}\) 均为等差数列，则 \(M\) 中最多有 1 个元素；
\item 若 \(\{a_n\}\) 与 \(\{b_n\}\) 均为等比数列，则 \(M\) 中最多有 2 个元素；
\item 若 \(\{a_n\}\) 为等差数列，\(\{b_n\}\) 为等比数列，则 \(M\) 中最多有 3 个元素；
\item 若 \(\{a_n\}\) 为递增数列，\(\{b_n\}\) 为递减数列，则 \(M\) 中最多有 1 个元素.
\end{circlelist}

其中正确结论的序号是\fillinblank{}.
\end{problem}

\section{解答题}

\begin{problem}
在 \(\triangle ABC\) 中，内角 \(A\)，\(B\)，\(C\) 的对边分别为 \(a\)，\(b\)，\(c\)，\(\angle A\) 为钝角，\(a=7\)，\(\sin 2B=\frac{\sqrt{3}}{7}b\cos B\).

\begin{enumerate}
\item 求 \(\angle A\)；
\item 从下列三个条件中选择一个作为已知，使得 \(\triangle ABC\) 存在，求 \(\triangle ABC\) 的面积.
\begin{circlelist}
\item \(b=7\)；
\item \(\cos B=\frac{13}{14}\)；
\item \(c\sin A=\frac{5\sqrt{3}}{2}\).
\end{circlelist}
\end{enumerate}
\end{problem}

\begin{problem}
如图，在四棱锥 \(P-ABCD\) 中，\(BC\parallel AD\)，\(AB=BC=1\)，\(AD=3\)，点 \(E\) 在 \(AD\) 上，且 \(PE\perp AD\)，\(PE=DE=2\).

\begin{enumerate}
\item 若 \(F\) 为线段 \(PE\) 中点，求证：\(BF\parallel \text{平面 }PCD\)；
\item 若 \(AB\perp \text{平面 }PAD\)，求平面 \(PAB\) 与平面 \(PCD\) 夹角的余弦值.
\end{enumerate}

\bitmapfigure[max width=5.1cm]{img/2024/beijing/q17_fig1.png}
\end{problem}

\begin{problem}
某保险公司为了了解该公司某种保险产品的索赔情况，从合同保险期限届满的保单中随机抽取 1000 份，记录并整理这些保单的索赔情况，获得数据如下表：

\begin{center}
\begin{tabular}{|c|c|c|c|c|c|}
\hline
索赔次数 & 0 & 1 & 2 & 3 & 4 \\
\hline
保单数 & 800 & 100 & 60 & 30 & 10 \\
\hline
\end{tabular}
\end{center}

假设：一份保单的保费为 0.4 万元；前 3 次索赔时，保险公司每次赔偿 0.8 万元；第四次索赔时，保险公司赔偿 0.6 万元. 假设不同保单的索赔次数相互独立. 用频率估计概率.

\begin{enumerate}
\item 估计一份保单索赔次数不少于 2 的概率；
\item 一份保单的毛利润定义为这份保单的保费与赔偿总金额之差.
    \begin{enumerate}
    \item 记 \(X\) 为一份保单的毛利润，估计 \(E(X)\)；
    \item 如果无索赔的保单的保费减少 4\%，有索赔的保单的保费增加 20\%，试比较这种情况下一份保单毛利润的数学期望估计值与（1）中 \(E(X)\) 估计值的大小（结论不要求证明）.
    \end{enumerate}
\end{enumerate}
\end{problem}

\begin{problem}
已知椭圆 \(E:\frac{x^2}{a^2}+\frac{y^2}{b^2}=1\)（\(a>b>0\)），以椭圆 \(E\) 的焦点和短轴端点为顶点的四边形是边长为 2 的正方形. 过点 \((0,t)\)（\(t>\sqrt{2}\)）且斜率存在的直线与椭圆 \(E\) 交于不同的两点 \(A\)，\(B\)，过点 \(A\) 和 \(C(0,1)\) 的直线 \(AC\) 与椭圆 \(E\) 的另一个交点为 \(D\).

\begin{enumerate}
\item 求椭圆 \(E\) 的方程及离心率；
\item 若直线 \(BD\) 的斜率为 0，求 \(t\) 的值.
\end{enumerate}
\end{problem}

\begin{problem}
设函数 \(f(x)=x+k\ln(1+x)\)（\(k\ne 0\)），直线 \(l\) 是曲线 \(y=f(x)\) 在点 \((t,f(t))\)（\(t>0\)）处的切线.

\begin{enumerate}
\item 当 \(k=-1\) 时，求 \(f(x)\) 的单调区间；
\item 求证：\(l\) 不经过点 \((0,0)\)；
\item 当 \(k=1\) 时，设点 \(A(t,f(t))\)（\(t>0\)），\(C(0,f(t))\)，\(O(0,0)\)，\(B\) 为 \(l\) 与 \(y\) 轴的交点，\(S_{\triangle ACO}\) 与 \(S_{\triangle ABO}\) 分别表示 \(\triangle ACO\) 与 \(\triangle ABO\) 的面积. 是否存在点 \(A\) 使得 \(2S_{\triangle ACO}=15S_{\triangle ABO}\) 成立？若存在，这样的点 \(A\) 有几个？
\end{enumerate}

（参考数据：\(1.09<\ln 3<1.10\)，\(1.60<\ln 5<1.61\)，\(1.94<\ln 7<1.95\)）
\end{problem}

\begin{problem}
已知集合 \(M=\left\{(i,j,k,w)\middle|\begin{gathered}i\in\{1,2\},\ j\in\{3,4\},\ k\in\{5,6\},\\ w\in\{7,8\},\ \text{且 }i+j+k+w\ \text{为偶数}\end{gathered}\right\}\).

给定数列 \(A:a_1\)，\(a_2\)，\(\cdots\)，\(a_8\)，和序列 \(\Omega:T_1\)，\(T_2\)，\(\cdots\)，\(T_s\)，其中 \(T_t=(i_t,j_t,k_t,w_t)\in M\)（\(t=1\)，\(2\)，\(\cdots\)，\(s\)）. 对数列 \(A\) 进行如下变换：将 \(A\) 的第 \(i_1\)，\(j_1\)，\(k_1\)，\(w_1\) 项均加 1，其余项不变，得到的数列记作 \(T_1(A)\)；将 \(T_1(A)\) 的第 \(i_2\)，\(j_2\)，\(k_2\)，\(w_2\) 项均加 1，其余项不变，得到数列记作 \(T_2T_1(A)\)；……；以此类推，得到 \(T_s\cdots T_2T_1(A)\)，简记为 \(\Omega(A)\).

\begin{enumerate}
\item 给定数列 \(A:1\)，\(3\)，\(2\)，\(4\)，\(6\)，\(3\)，\(1\)，\(9\) 和序列 \(\Omega:(1,3,5,7)\)，\((2,4,6,8)\)，\((1,3,5,7)\)，写出 \(\Omega(A)\)；
\item 是否存在序列 \(\Omega\)，使得 \(\Omega(A)\) 为 \(a_1+2\)，\(a_2+6\)，\(a_3+4\)，\(a_4+2\)，\(a_5+8\)，\(a_6+2\)，\(a_7+4\)，\(a_8+4\)？若存在，写出一个符合条件的 \(\Omega\)；若不存在，请说明理由；
\item 若数列 \(A\) 的各项均为正整数，且 \(a_1+a_3+a_5+a_7\) 为偶数，求证：“存在序列 \(\Omega\)，使得 \(\Omega(A)\) 的各项都相等”的充要条件为“\(a_1+a_2=a_3+a_4=a_5+a_6=a_7+a_8\)”.
\end{enumerate}
\end{problem}

